\subsubsection{Diagnostic Classification and MCI Stage Discrimination}
Image-based classifiers used a DenseNet121 backbone~\cite{huang2017densely} to evaluate three cross-sectional tasks: AD/CN; EMCI+LMCI/CN; and EMCI/LMCI. Positive classes were AD, EMCI+LMCI and LMCI, respectively. The task-specific cohorts contained 470, 741 and 436 participants, excluding MCI-U. Seven imaging conditions used MRI alone or MRI combined with real/synthetic FDG, AV45 or both tracers. Two further conditions were Clinical-only (an 18-parameter linear classifier) and MRI+Clinical. Both clinical baselines used the same eight clinical variables as the conditioning inputs; their information content is diagnosis-related and affected by historical group-conditioned filling.

All conditions used the same task-specific subject-level stratified three-fold outer splits (seed 42). Within each outer training set, a fixed stratified 15\% validation subset selected the checkpoint by validation ACC; the outer test was evaluated once after reloading that checkpoint. Imaging models used Adam (learning rate $10^{-4}$), batch size 5, up to 200 epochs and validation every two epochs. Clinical and combined-model settings are retained separately in the reproducibility records rather than assumed identical to every imaging model. Thresholded metrics use 0.5; AUC uses continuous positive-class scores. We report ACC, AUC, SEN, SPE and balanced ACC as equally weighted means across the three folds with sample s.d.

Bootstrap CIs use 5,000 paired resamples within each fold-by-class stratum (seed 20261102); each condition shares participant multiplicities within a resample. Metrics are recalculated by fold and averaged equally, rather than pooling differently trained model scores. Paired differences use the same resamples. The 95\% percentile intervals are marginal and unadjusted for simultaneous coverage; they condition on the trained models, class proportions and fixed splits and do not estimate retraining variability or eliminate dependence from overlapping training sets. DeLong tests compare matched predictions separately within each outer fold. The 18 original imaging comparisons and 36 additional clinical comparisons form separate Holm families; raw $p$ values are retained. No overall three-fold DeLong $p$ value is reported. ROC curves average fold-specific TPR over an FPR grid with step 0.005; pointwise 95\% bands use 2,000 fold-by-class resamples. The legend AUC is the mean empirical fold AUC, not pooled AUC or an integral of the finite grid.

Similar observed Real/Syn performance or a non-significant difference does not establish biological fidelity, equivalence or clinical interchangeability. Clinical-only and MRI+Clinical comparisons assess whether the reported synthetic-image classifier adds information beyond the provided clinical fields; they do not resolve label circularity introduced by diagnostic-group-conditioned simulated inputs.

Demographic sensitivity analyses separately stratified recorded age at 75 years and sex (F/M), without crossing the strata or changing thresholds/checkpoints. Each task, fold, subgroup and true-class stratum was resampled 5,000 times (seed 20261103); seven imaging inputs shared resample multiplicities, and fold metrics were averaged equally. These unadjusted marginal subgroup CIs were not used to infer interactions. Recorded ages could not be verified as MRI-visit ages because the available demographic export lacked visit dates. Age and sex were not mutually adjusted, and clinical-score or residual within-group confounding was not eliminated.
